\documentclass[12pt]{article}
\usepackage{amsmath}
\usepackage{amssymb}
\usepackage{fullpage}
\usepackage{overcite}
\renewcommand{\familydefault}{\sfdefault}

% separate paragraphs with vertical space rather than indentation
\setlength{\parskip}{8pt plus 2pt minus 1pt}
\setlength{\parindent}{0pt}

% antisymmetrized two-electron integral <pq||rs>
\newcommand{\adb}[2]{\langle #1 \| #2 \rangle}
\newcommand{\so}{\text{all}}
\newcommand{\occ}{\text{occ}}
\newcommand{\vir}{\text{vir}}

\begin{document}

\begin{center}
{\bf \large Correlated Derivatives with Orbital Response} \\[2pt]
T.~Daniel Crawford \\
{\em Virginia Tech, Blacksburg, Virginia} \\
\end{center}

Derivatives of electronic energies with respect to external perturbations can
be elegantly formulated in terms of electron densities regardless of method or
level of theory.\cite{Yamaguchi94}  In these notes we will explore derivatives
of correlation energies with respect to {\em real} perturbations, especially the
handling of the response of the underlying molecular orbitals (MOs).  While we
will compare many of the expressions we derive here to literature references,
especially for coupled cluster theory, the results are sufficiently general so
as to apply to any correlated method.

\section{First Derivatives}

The first derivative of the correlation energy with respect to a {\em real}
perturbation (e.g., a nuclear-coordinate displacement or a static external electric
field) may be written as:\cite{Gauss91:UHF}
\begin{equation}
\frac{\partial E}{\partial x}
= \sum_{pq} D_{pq}\,\frac{\partial f_{pq}}{\partial x}
+ \sum_{pqrs} \Gamma_{pqrs}\,\frac{\partial \adb{pq}{rs}}{\partial x},
\label{eq:dE}
\end{equation}
where $D_{pq}$ and $\Gamma_{pqrs}$ are the (unrelaxed) one- and two-particle
density matrices in the spin-orbital basis and $\Gamma_{pqrs}$ is antisymmetric.

\noindent
The antisymmetrized two-electron integrals recur in the pair combination
$A_{pqrs}\equiv\adb{pq}{rs}+\adb{ps}{rq}$.  The derivatives of the Fock matrix
elements and of the two-electron integrals are
\begin{align}
\frac{\partial f_{pq}}{\partial x}
&= f^{(x)}_{pq} + U^x_{qp} f_{qq} + U^x_{pq} f_{pp}
+ \sum_{t}^{\so}\sum_{m}^{\occ} U^x_{tm}A_{ptqm},
\label{eq:dfpq}\\[4pt]
\frac{\partial \adb{pq}{rs}}{\partial x}
&= \adb{pq}{rs}^{(x)}
+ \sum_{t}^{\so}\left(U^x_{tp}\adb{tq}{rs} + U^x_{tq}\adb{pt}{rs}
+ U^x_{tr}\adb{pq}{ts} + U^x_{ts}\adb{pq}{rt}\right),
\label{eq:deri}
\end{align}
where $f^{(x)}_{pq}$ and $\adb{pq}{rs}^{(x)}$ are the ``skeleton'' derivative
integrals (AO-basis derivatives transformed into the unperturbed MO basis) and
the $U^x_{pq}$ are the coupled-perturbed Hartree--Fock (CPHF) coefficients.

\noindent
Inserting \eqref{eq:dfpq} and \eqref{eq:deri} into \eqref{eq:dE} and collecting
the terms with and without $U^{x}$,
\begin{align}
\frac{\partial E}{\partial x}
&= \sum_{pq} D_{pq} f^{(x)}_{pq} + \sum_{pqrs}\Gamma_{pqrs}\adb{pq}{rs}^{(x)}
+ \sum_{pq} D_{pq}\left[U^x_{qp} f_{qq} + U^x_{pq} f_{pp}\right]
\nonumber\\
&\quad
+ \sum_{pq}\sum_{t}^{\so}\sum_{m}^{\occ} D_{pq}\,U^x_{tm}A_{ptqm}
\nonumber\\
&\quad
+ \sum_{pqrs}\sum_{t}^{\so}\Gamma_{pqrs}\left(U^x_{tp}\adb{tq}{rs}
+ U^x_{tq}\adb{pt}{rs} + U^x_{tr}\adb{pq}{ts} + U^x_{ts}\adb{pq}{rt}\right).
\label{eq:collect}
\end{align}

\noindent
Re-indexing every $U^{x}$-dependent term so that the CPHF coefficient carries
the labels $pq$ (dummies relabeled $r,s,t$) gives
\begin{align}
\frac{\partial E}{\partial x}
&= \sum_{pq} D_{pq} f^{(x)}_{pq} + \sum_{pqrs}\Gamma_{pqrs}\adb{pq}{rs}^{(x)}
+ \sum_{pq} U^x_{pq} f_{pp}\left(D_{pq}+D_{qp}\right)
\nonumber\\
&\quad
+ \sum_{p}^{\so}\sum_{q} U^x_{pq}\,\delta_{q,\occ}\sum_{rs}
D_{rs}A_{rpsq}
\nonumber\\
&\quad
+ \sum_{pq} U^x_{pq}\sum_{rst}\left(
\adb{pr}{st}\,\Gamma_{qrst}
+ \adb{rp}{st}\,\Gamma_{rqst}
+ \adb{rs}{pt}\,\Gamma_{rsqt}
+ \adb{rs}{tp}\,\Gamma_{rstq}\right).
\label{eq:reindex}
\end{align}

\noindent
The four integral-rotation terms are pairwise equal through the antisymmetry of
$\Gamma$ (terms 1--2, since $\adb{rp}{st}=-\adb{pr}{st}$ and
$\Gamma_{rqst}=-\Gamma_{qrst}$; terms 3--4 likewise), and equal to one another
once $\Gamma$ carries the two-particle bra--ket symmetry
$\Gamma_{pqrs}=\Gamma_{rspq}$; they therefore combine to
$4\sum_{rst}\adb{pr}{st}\Gamma_{qrst}$.  Finally,
\begin{equation}
\frac{\partial E}{\partial x}
= \sum_{pq} D_{pq} f^{(x)}_{pq}
+ \sum_{pqrs}\Gamma_{pqrs}\adb{pq}{rs}^{(x)}
- 2\sum_{pq} U^x_{pq}\, I'_{pq},
\label{eq:dE-Iprime}
\end{equation}
where
\begin{equation}
I'_{pq} = -\tfrac{1}{2}\left[
f_{pp}\left(D_{pq}+D_{qp}\right)
+ \delta_{q,\occ}\sum_{rs} D_{rs}A_{rpsq}
+ 4\sum_{rst}\adb{pr}{st}\Gamma_{qrst}
\right].
\label{eq:Iprime}
\end{equation}
The expression for $I'_{pq}$ assumes $\Gamma_{pqrs}$ is antisymmetric but
$D_{pq}$ is not necessarily symmetric.  The prefactor of $-\tfrac{1}{2}$ is
standard for most CC/CI gradient papers and leads to simpler later expressions.
Equations \eqref{eq:dE-Iprime} and \eqref{eq:Iprime} match Eqs.~42 and 43 of
Ref.~\citenum{Gauss91:UHF}.

\section{Simplifying the Orbital-Gradient Term}

Separate the orbital-gradient term into MO subspaces:
\begin{equation}
-2\sum_{pq} U^x_{pq} I'_{pq}
= -2\sum_{ij} U^x_{ij} I'_{ij}
- 2\sum_{ab} U^x_{ab} I'_{ab}
- 2\sum_{ai} U^x_{ai} I'_{ai}
- 2\sum_{ia} U^x_{ia} I'_{ia}.
\end{equation}
The derivative of the MO orthonormality conditions gives
\begin{equation}
U^x_{pq} + U^x_{qp} + S^{(x)}_{pq} = 0,
\label{eq:ortho}
\end{equation}
so the last term may be rewritten as
$-2\sum_{ia}U^x_{ia}I'_{ia} = +2\sum_{ia}\left(U^x_{ai}+S^{(x)}_{ai}\right)I'_{ia}$,
and the orbital-gradient term becomes
\begin{equation}
-2\sum_{pq} U^x_{pq} I'_{pq}
= -2\sum_{ij} U^x_{ij} I'_{ij}
- 2\sum_{ab} U^x_{ab} I'_{ab}
+ 2\sum_{ai} U^x_{ai}\left(I'_{ia}-I'_{ai}\right)
+ 2\sum_{ai} S^{(x)}_{ai} I'_{ia}.
\label{eq:orbgrad}
\end{equation}

\noindent
Using the {\em non-canonical perturbed-orbital
conditions}\cite{Handy85:MP2Hessians}
\begin{equation}
U^x_{ij} = -\tfrac{1}{2}S^{(x)}_{ij}, \qquad U^x_{ab} = -\tfrac{1}{2}S^{(x)}_{ab},
\label{eq:noncanon}
\end{equation}
yields
\begin{equation}
-2\sum_{pq} U^x_{pq} I'_{pq}
= \sum_{ij} S^{(x)}_{ij} I'_{ij}
+ \sum_{ab} S^{(x)}_{ab} I'_{ab}
+ 2\sum_{ai} U^x_{ai}\left(I'_{ia}-I'_{ai}\right)
+ 2\sum_{ai} S^{(x)}_{ai} I'_{ia}.
\end{equation}
All terms involving overlap derivatives may be collected together after
splitting the last one, $2\sum_{ai}S^{(x)}_{ai}I'_{ia} = \sum_{ai}S^{(x)}_{ai}I'_{ia}
+ \sum_{ia}S^{(x)}_{ia}I'_{ia}$:
\begin{equation}
\sum_{ij} S^{(x)}_{ij} I'_{ij}
+ \sum_{ab} S^{(x)}_{ab} I'_{ab}
+ \sum_{ia} S^{(x)}_{ia} I'_{ia}
+ \sum_{ai} S^{(x)}_{ai} I'_{ia}
= \sum_{pq} S^{(x)}_{pq} I''_{pq},
\end{equation}
where
\begin{equation}
I''_{pq} =
\begin{cases}
I'_{qp}, & (p,q)=(a,i),\\
I'_{pq}, & \text{otherwise.}
\end{cases}
\label{eq:Idouble}
\end{equation}

\noindent
The total gradient may therefore be written as
\begin{equation}
\frac{\partial E}{\partial x}
= \sum_{pq} D_{pq} f^{(x)}_{pq}
+ \sum_{pqrs}\Gamma_{pqrs}\adb{pq}{rs}^{(x)}
+ 2\sum_{a}^{\vir}\sum_{i}^{\occ} U^x_{ai} X_{ai}
+ \sum_{pq} S^{(x)}_{pq} I''_{pq},
\qquad X_{ai} \equiv I'_{ia} - I'_{ai}
\label{eq:grad-U}
\end{equation}
(Compare to Eq.~47 of Ref.~\citenum{Gauss91:UHF}.)

\section{CPHF Equations and the Z-Vector}

The first-order CPHF equations for $U^x_{ai}$ are
\begin{equation}
-f^{(x)}_{ai} + S^{(x)}_{ai}\varepsilon_i
+ \tfrac{1}{2}\sum_{jk} S^{(x)}_{jk}A_{ajik}
= \sum_{b}^{\vir}\sum_{j}^{\occ} U^x_{bj}
\left[(\varepsilon_a-\varepsilon_i)\delta_{ab}\delta_{ij} + A_{abij}\right].
\end{equation}
Defining
\begin{align}
G_{ai,bj} &\equiv (\varepsilon_a-\varepsilon_i)\delta_{ab}\delta_{ij} + A_{abij},\\
B^x_{ai} &\equiv -f^{(x)}_{ai} + S^{(x)}_{ai}\varepsilon_i
+ \tfrac{1}{2}\sum_{jk} S^{(x)}_{jk}A_{ajik},
\end{align}
the CPHF equations read $\sum_{bj} G_{ai,bj} U^x_{bj} = B^x_{ai}$, or
$\mathbf{G}\,\mathbf{U}^{x} = \mathbf{B}^{x}$.

\noindent
The orbital contribution to the gradient is
$\dfrac{\partial E}{\partial x} \leftarrow 2\sum_{ai} U^x_{ai} X_{ai}
= 2\,\mathbf{X}^{T}\mathbf{U}^{x}$.
Following Handy and Schaefer,\cite{Handy84}
$\mathbf{U}^{x} = \mathbf{G}^{-1}\mathbf{B}^{x}$ so that
$\dfrac{\partial E}{\partial x} \leftarrow 2\,\mathbf{X}^{T}\mathbf{G}^{-1}\mathbf{B}^{x}$
(cf.\ Eq.~51 of Ref.~\citenum{Gauss91:UHF}, same except for sign).  Defining
\begin{equation}
\mathbf{Z}^{T} \equiv \mathbf{X}^{T}\mathbf{G}^{-1}
\;\Longrightarrow\;
\mathbf{G}^{T}\mathbf{Z} = \mathbf{X}, \quad\text{i.e.}\quad
\sum_{bj} G_{ai,bj}^{T} Z_{bj} = X_{ai},
\end{equation}
gives the interchanged (Z-vector) form
\begin{equation}
\frac{\partial E}{\partial x} \leftarrow 2\,\mathbf{Z}^{T}\mathbf{B}^{x}
= 2\sum_{ai} Z_{ai} B^x_{ai}.
\end{equation}
Expanding $B^x_{ai}$,
\begin{equation}
\frac{\partial E}{\partial x} \leftarrow
-2\sum_{ai} Z_{ai} f^{(x)}_{ai}
+ 2\sum_{ai} Z_{ai} S^{(x)}_{ai}\varepsilon_i
+ \sum_{ai} Z_{ai}\sum_{jk} S^{(x)}_{jk}A_{ajik},
\end{equation}
and, re-indexing the third term and splitting the first two,
\begin{align}
\frac{\partial E}{\partial x} \leftarrow
&-\sum_{ai} Z_{ai} f^{(x)}_{ai} - \sum_{ia} Z_{ai} f^{(x)}_{ia}
+ \sum_{ai} Z_{ai} S^{(x)}_{ai}\varepsilon_i + \sum_{ai} Z_{ai} S^{(x)}_{ai}\varepsilon_i
\nonumber\\
&+ \sum_{ij} S^{(x)}_{ij}\sum_{ak} Z_{ak}A_{aikj}.
\end{align}

\section{Total Gradient}

Collecting everything, the total gradient becomes
\begin{equation}
\frac{\partial E}{\partial x}
= \sum_{pq} \tilde{D}_{pq} f^{(x)}_{pq}
+ \sum_{pqrs}\Gamma_{pqrs}\adb{pq}{rs}^{(x)}
+ \sum_{pq} I_{pq} S^{(x)}_{pq},
\label{eq:grad-final}
\end{equation}
where
\begin{align}
\tilde{D}_{ai} &= D_{ai} - Z_{ai}, & \tilde{D}_{ia} &= D_{ia} - Z_{ai}, \\
\tilde{D}_{ij} &= D_{ij}, & \tilde{D}_{ab} &= D_{ab}, \\
I_{ai} &= I''_{ai} + Z_{ai}\varepsilon_i, & I_{ia} &= I''_{ia} + Z_{ai}\varepsilon_i, \\
I_{ij} &= I''_{ij} + \sum_{a}^{\vir}\sum_{k}^{\occ} Z_{ak}A_{aikj}, & I_{ab} &= I''_{ab}.
\end{align}

\noindent
{\em Note:} To compare the above to Eqs.~52--56 of Ref.~\citenum{Gauss91:UHF},
be aware that the orbital Z-vector equations used here differ by a sign from
those in Eq.~51 of that paper.  Hence the signs on the Z-vector contributions
are reversed; otherwise everything matches.

\section{Canonical Perturbed Orbitals}

The reference derivation of \S2 uses the non-canonical perturbed-orbital
conditions \eqref{eq:noncanon}, $U^x_{ij}=-\tfrac12 S^{(x)}_{ij}$ and
$U^x_{ab}=-\tfrac12 S^{(x)}_{ab}$.  These are permissible only when the correlation
energy is invariant to rotations within the occupied space and within the
virtual space.  Equation~\eqref{eq:Iprime} is {\em not} manifestly symmetric in
$(i,j)$ or $(a,b)$, since neither the $f_{pp}$ prefactor ($f_{ii}$ vs.\ $f_{jj}$) nor
the three-index contraction ($\adb{ir}{st}\Gamma_{jrst}$ vs.\
$\adb{jr}{st}\Gamma_{irst}$) is symmetric under the swap.  What matters instead is
its {\em antisymmetric} part, which enters in exactly the way
$X_{ai}=I'_{ia}-I'_{ai}$ did for the occ-vir block in \S2.  Applying the
orthonormality condition \eqref{eq:ortho} to the occupied--occupied and
virtual--virtual blocks of the subspace split of $-2\sum_{pq}U^x_{pq}I'_{pq}$
(relabeling $i\leftrightarrow j$, resp.\ $a\leftrightarrow b$, in half of each
sum and eliminating $U^x_{ji}=-S^{(x)}_{ij}-U^x_{ij}$) gives
\begin{align}
-2\sum_{ij}U^x_{ij}I'_{ij}
&= -\sum_{ij}U^x_{ij}\left(I'_{ij}-I'_{ji}\right) + \sum_{ij}S^{(x)}_{ij}I'_{ij},
\label{eq:orbgrad-oo}\\
-2\sum_{ab}U^x_{ab}I'_{ab}
&= -\sum_{ab}U^x_{ab}\left(I'_{ab}-I'_{ba}\right) + \sum_{ab}S^{(x)}_{ab}I'_{ab},
\label{eq:orbgrad-vv}
\end{align}
so the antisymmetric orbital-rotation gradients $I'_{ij}-I'_{ji}$ and
$I'_{ab}-I'_{ba}$ now appear explicitly.  Combining these with the occ-vir block
already reduced in Eq.~\eqref{eq:orbgrad}, and collapsing every overlap-derivative
term into the single sum $\sum_{pq}S^{(x)}_{pq}I''_{pq}$ (with $I''$ as defined in
Eq.~\eqref{eq:Idouble}, since the occ-occ, vir-vir, and occ-vir overlap terms are identical to
those of \S2), the orbital-gradient term becomes
\begin{equation}
-2\sum_{pq}U^x_{pq}I'_{pq}
= -\sum_{ij}U^x_{ij} X_{ij}
- \sum_{ab}U^x_{ab} X_{ab}
+ 2\sum_{ai}U^x_{ai} X_{ai} 
+ \sum_{pq}S^{(x)}_{pq}I''_{pq},
\label{eq:orbgrad-canon-full}
\end{equation}
where $X_{ij}\equiv I'_{ij}-I'_{ji}$ and $X_{ab}\equiv I'_{ab}-I'_{ba}$ are the
occupied--occupied and virtual--virtual analogs of $X_{ai}=I'_{ia}-I'_{ai}$ in
Eq.~\eqref{eq:grad-U}.  For quantum chemical methods for which the energy is invariant to occ--occ and
virt--virt rotations (most methods), $X_{ij}=X_{ab}=0$ (a property of the 
densities $D,\Gamma$, not a manifest symmetry of \eqref{eq:Iprime}).  The
$U^x_{ij}X_{ij}$ and $U^x_{ab}X_{ab}$ terms then vanish for any choice of the
occupied and virtual rotations, so those rotations may be set to their
orthonormality values $-\tfrac12 S^{(x)}$ and Eq.~\eqref{eq:grad-U} is recovered.
For a method whose energy is not invariant to occupied--occupied or
virtual--virtual rotations (so that $X_{ij}\neq0$ and $X_{ab}\neq0$, as for the
canonical-MO (T) estimate\cite{Lee91,Hald03:gradients}), or if core orbitals are
frozen (because the correlated energy is not invariant to rotations between core
MOs and active occupied MOs), then $U^x_{ij}$ and $U^x_{ab}$ must be explicitly
computed using the {\em canonical} perturbed-orbital conditions
\begin{equation}
\frac{\partial f_{ij}}{\partial x}=0 \quad(i\neq j),\qquad
\frac{\partial f_{ab}}{\partial x}=0 \quad(a\neq b),
\label{eq:canon}
\end{equation}
i.e.\ the perturbed occupied (virtual) orbitals are kept canonical, diagonalizing
the occupied (virtual) block of the perturbed Fock matrix.

\noindent
The first of Eqs.~\eqref{eq:canon} along with $f_{pp}=\varepsilon_p$
for canonical-MOs and orthonormality $U^x_{ji}=-S^{(x)}_{ij}-U^x_{ij}$,
determines $U^x_{ij}$ ($i\neq j$) fully:
\begin{equation}
U^x_{ij}
= \frac{1}{\varepsilon_i-\varepsilon_j}\left[
-f^{(x)}_{ij} + S^{(x)}_{ij}\varepsilon_j
+ \tfrac12\sum_{nm}^{\occ} S^{(x)}_{nm}\,A_{injm}
- \sum_{c}^{\vir}\sum_{m}^{\occ} U^x_{cm}\,A_{icjm}
\right], \qquad i\neq j,
\label{eq:canon-oo}
\end{equation}
and, analogously, the virtual block
\begin{equation}
U^x_{ab}
= \frac{1}{\varepsilon_a-\varepsilon_b}\left[
-f^{(x)}_{ab} + S^{(x)}_{ab}\varepsilon_b
+ \tfrac12\sum_{nm}^{\occ} S^{(x)}_{nm}\,A_{anbm}
- \sum_{c}^{\vir}\sum_{m}^{\occ} U^x_{cm}\,A_{acbm}
\right], \qquad a\neq b.
\label{eq:canon-vv}
\end{equation}
The diagonal elements remain $U^x_{ii}=-\tfrac12 S^{(x)}_{ii}$,
$U^x_{aa}=-\tfrac12 S^{(x)}_{aa}$ from orthonormality.  

While equations \eqref{eq:canon-oo}--\eqref{eq:canon-vv} provide an alternative
route to accounting for the response of the MOs to the external (real)
perturbation, they also introduce a numerical sensitivity in the case of
(near-)degenerate occupied or virtual orbitals
through the energy-denominators
$\varepsilon_i-\varepsilon_j$ and
$\varepsilon_a-\varepsilon_b$.\cite{Handy85:MP2Hessians}

\noindent
Assembling, the total gradient in the canonical-orbital formulation is the analog
of Eq.~\eqref{eq:grad-U} with the occupied--occupied and virtual--virtual
orbital-response terms retained,
\begin{equation}
\frac{\partial E}{\partial x}
= \sum_{pq} D_{pq} f^{(x)}_{pq}
+ \sum_{pqrs}\Gamma_{pqrs}\adb{pq}{rs}^{(x)}
- \sum_{ij} U^x_{ij} X_{ij}
- \sum_{ab} U^x_{ab} X_{ab}
+ 2\sum_{ai} U^x_{ai} X_{ai}
+ \sum_{pq} S^{(x)}_{pq} I''_{pq},
\label{eq:grad-canon}
\end{equation}
with $U^x_{ai}$ obtained from the CPHF/Z-vector of \S3 and $U^x_{ij}$, $U^x_{ab}$
from Eqs.~\eqref{eq:canon-oo} and \eqref{eq:canon-vv}.  If the correlation
method is invariant to occ-occ and vir-vir MO rotations, the two new terms
vanish ($X_{ij}=X_{ab}=0$) and Eq.~\eqref{eq:grad-U} is recovered.  Note that $D_{pq}$ is still the unrelaxed
one-electron density in this expression, comparable to Eq.~\eqref{eq:grad-U} where we used the non-canonical perturbed orbital
approach.

\section{CPHF and Z-Vector for Canonical-Perturbed MOs}

In Eq.~\eqref{eq:grad-canon} the occ-vir orbital-response term $2\sum_{ai}U^x_{ai}X_{ai}$
is handled exactly as in \S3: the occ-vir coefficients solve
$\mathbf{G}\,\mathbf{U}^{x}=\mathbf{B}^{x}$, and the interchange
$\mathbf{G}^{T}\mathbf{Z}=\mathbf{X}$ trades the per-perturbation solve for a single
Z-vector.  What is new for the canonical perturbed MO case is the occupied--occupied and
virtual--virtual terms $-\sum_{ij}U^x_{ij}X_{ij}$ and $-\sum_{ab}U^x_{ab}X_{ab}$,
whose rotations are given {\em explicitly} by Eqs.~\eqref{eq:canon-oo} and
\eqref{eq:canon-vv} rather than by an iterative solve.

\noindent
Substituting Eq.~\eqref{eq:canon-oo} and introducing the occupied--occupied
{\em dependent-pair}
\begin{equation}
P_{ij} \equiv \frac{X_{ij}}{\varepsilon_i-\varepsilon_j}
= \frac{I'_{ij}-I'_{ji}}{\varepsilon_i-\varepsilon_j},
\label{eq:Poo}
\end{equation}
the occupied block becomes
\begin{equation}
-\sum_{ij}U^x_{ij}X_{ij}
= \sum_{ij}P_{ij}f^{(x)}_{ij}
- \sum_{ij}P_{ij}S^{(x)}_{ij}\varepsilon_j
- \tfrac12\sum_{ij}\sum_{nm}^{\occ}P_{ij}S^{(x)}_{nm}A_{injm}
+ \sum_{ij}\sum_{c}^{\vir}\sum_{m}^{\occ}P_{ij}U^x_{cm}A_{icjm}.
\label{eq:oo-sub}
\end{equation}
The four terms have distinct
destinations: the first adds $P_{ij}$ to the effective
one-particle density that contracts the skeleton Fock derivative; the second and
third are overlap-derivative contributions that join the energy-weighted density
$I''$; and the fourth couples the occ-vir response $U^x_{cm}$, so it augments the
Z-vector right-hand side.

\noindent
The virtual block follows identically, with
$P_{ab}\equiv X_{ab}/(\varepsilon_a-\varepsilon_b)$,
\begin{equation}
-\sum_{ab}U^x_{ab}X_{ab}
= \sum_{ab}P_{ab}f^{(x)}_{ab}
- \sum_{ab}P_{ab}S^{(x)}_{ab}\varepsilon_b
- \tfrac12\sum_{ab}\sum_{nm}^{\occ}P_{ab}S^{(x)}_{nm}A_{anbm}
+ \sum_{ab}\sum_{c}^{\vir}\sum_{m}^{\occ}P_{ab}U^x_{cm}A_{acbm}.
\label{eq:vv-sub}
\end{equation}

Collecting the terms of Eqs.~\eqref{eq:oo-sub} and \eqref{eq:vv-sub} sends each
of them to one of the three density-like structures already present in the
gradient, Eq.~\eqref{eq:grad-canon}.

The leading terms, $P_{ij}f^{(x)}_{ij}$ and $P_{ab}f^{(x)}_{ab}$, have the same
shape as the skeleton Fock term $\sum_{pq}D_{pq}f^{(x)}_{pq}$, so the
dependent-pairs are absorbed by augmenting the one-particle density,
\begin{equation}
\tilde{D}_{ij} = D_{ij} + P_{ij}, \qquad \tilde{D}_{ab} = D_{ab} + P_{ab}.
\label{eq:Dtilde}
\end{equation}

The terms carrying an overlap derivative $S^{(x)}$ match the energy-weighted density
term $\sum_{pq}S^{(x)}_{pq}I''_{pq}$, so they are absorbed into $I''$,
\begin{equation}
\tilde{I}''_{ij} = I''_{ij} - P_{ij}\varepsilon_j
- \tfrac12\sum_{kl}^{\occ}P_{kl}A_{kilj}
- \tfrac12\sum_{cd}^{\vir}P_{cd}A_{cidj},
\qquad
\tilde{I}''_{ab} = I''_{ab} - P_{ab}\varepsilon_b
\label{eq:Idouble-tilde}
\end{equation}
(both $A$ sums in the occupied--occupied block come from the third terms of
Eqs.~\eqref{eq:oo-sub} and \eqref{eq:vv-sub}, which multiply an
occupied--occupied overlap $S^{(x)}_{nm}$).

The final terms carry an occ-vir rotation coefficient $U^x_{cm}$ ($c$ virtual, $m$
occupied).  Unlike the occ-occ and vir-vir rotations, which the canonical conditions fix
in closed form through Eqs.~\eqref{eq:canon-oo} and \eqref{eq:canon-vv}, the
occ-vir
rotations remain the unknowns of the CPHF problem, so these terms cannot be
evaluated directly and must instead be carried through the Z-vector.  Because
$U^x_{cm}$ sits in the same occ-vir block as the orbital term
$2\sum_{ai}U^x_{ai}X_{ai}$, renaming the dummy label $(c,m)\to(a,i)$ lets the two
be added under a single sum $\sum_{ai}U^x_{ai}(\cdots)$; factoring the common $2$
then augments the Z-vector right-hand side $X_{ai}$ of \S3 to
\begin{equation}
\tilde{X}_{ai} = X_{ai}
+ \tfrac12\sum_{kl}^{\occ}P_{kl}\,A_{kali}
+ \tfrac12\sum_{de}^{\vir}P_{de}\,A_{daei}.
\label{eq:Xtilde}
\end{equation}
The interchange $\mathbf{G}^{T}\mathbf{Z}=\tilde{\mathbf{X}}$ and the expansion of
$B^x_{ai}$ then proceed exactly as in \S3, now with $\tilde{D}$, $\tilde{I}''$,
and $\tilde{X}$ in place of $D$, $I''$, and $X$.

\section{Computational Summary}

Both perturbed-orbital choices lead to the same compact gradient,
Eq.~\eqref{eq:grad-final}, repeated here,
\begin{equation}
\frac{\partial E}{\partial x}
= \sum_{pq} \tilde{D}_{pq} f^{(x)}_{pq}
+ \sum_{pqrs}\Gamma_{pqrs}\adb{pq}{rs}^{(x)}
+ \sum_{pq} I_{pq} S^{(x)}_{pq},
\tag{\ref{eq:grad-final}}
\end{equation}
and differ only in how the relaxed one-particle density $\tilde{D}$ and the
energy-weighted density $I$ are assembled from the unrelaxed densities $D$ and
$\Gamma$.

\subsection*{Non-Canonical Orbitals}
\begin{enumerate}
\item Build $I'_{pq}$ from Eq.~\eqref{eq:Iprime} and $I''_{pq}$ from
Eq.~\eqref{eq:Idouble}.
\item Solve the single Z-vector equation $\mathbf{G}^{T}\mathbf{Z}=\mathbf{X}$
with $X_{ai}=I'_{ia}-I'_{ai}$.
\item Assemble the densities,
\begin{equation*}
\tilde{D}_{ai}=D_{ai}-Z_{ai},\quad \tilde{D}_{ia}=D_{ia}-Z_{ai},\qquad
\tilde{D}_{ij}=D_{ij},\quad \tilde{D}_{ab}=D_{ab},
\end{equation*}
\begin{equation*}
I_{ai}=I''_{ai}+Z_{ai}\varepsilon_i,\quad
I_{ia}=I''_{ia}+Z_{ai}\varepsilon_i,\quad
I_{ij}=I''_{ij}+\sum_{a}^{\vir}\sum_{k}^{\occ} Z_{ak}A_{aikj},\quad
I_{ab}=I''_{ab}.
\end{equation*}
\end{enumerate}

\subsection*{Canonical Orbitals}
\begin{enumerate}
\item Build $I'_{pq}$ and $I''_{pq}$ as above.
\item Form the occupied--occupied and virtual--virtual dependent-pairs
(Eq.~\eqref{eq:Poo}),
\begin{equation*}
P_{ij}=\frac{I'_{ij}-I'_{ji}}{\varepsilon_i-\varepsilon_j},\qquad
P_{ab}=\frac{I'_{ab}-I'_{ba}}{\varepsilon_a-\varepsilon_b}.
\end{equation*}
\item Fold the dependent-pairs into the intermediates {\em before} the Z-vector
solve.  Augment the one-particle density,
\begin{equation*}
\tilde{D}_{ij}=D_{ij}+P_{ij},\qquad \tilde{D}_{ab}=D_{ab}+P_{ab};
\end{equation*}
the energy-weighted density,
\begin{equation*}
\tilde{I}''_{ij} = I''_{ij} - P_{ij}\varepsilon_j
- \tfrac12\sum_{kl}^{\occ}P_{kl}A_{kilj}
- \tfrac12\sum_{cd}^{\vir}P_{cd}A_{cidj},
\qquad
\tilde{I}''_{ab} = I''_{ab} - P_{ab}\varepsilon_b;
\end{equation*}
and the Z-vector right-hand side,
\begin{equation*}
\tilde{X}_{ai} = X_{ai}
+ \tfrac12\sum_{kl}^{\occ}P_{kl}A_{kali}
+ \tfrac12\sum_{de}^{\vir}P_{de}A_{daei}.
\end{equation*}
\item Solve the single Z-vector equation
$\mathbf{G}^{T}\mathbf{Z}=\tilde{\mathbf{X}}$.
\item Complete the assembly as in the non-canonical case, using the $\tilde{D}$
blocks from step~3 and $\tilde{I}''$ in place of $I''$,
\begin{equation*}
\tilde{D}_{ai}=D_{ai}-Z_{ai},\quad \tilde{D}_{ia}=D_{ia}-Z_{ai},
\end{equation*}
\begin{equation*}
I_{ai}=\tilde{I}''_{ai}+Z_{ai}\varepsilon_i,\quad
I_{ia}=\tilde{I}''_{ia}+Z_{ai}\varepsilon_i,\quad
I_{ij}=\tilde{I}''_{ij}+\sum_{a}^{\vir}\sum_{k}^{\occ} Z_{ak}A_{aikj},\quad
I_{ab}=\tilde{I}''_{ab}.
\end{equation*}
\end{enumerate}

For methods for which the energy is invariant to occupied--occupied and virtual--virtual
rotations, the dependent-pairs vanish ($X_{ij}=X_{ab}=0$, hence
$P_{ij}=P_{ab}=0$), $\tilde{I}''=I''$, and $\tilde{X}=X$; the canonical recipe
then reduces term by term to the non-canonical one.  The two routes diverge only
when we choose to use canonical perturbed MOs (as is typically the choice for
CCSD(T)).

\appendix

\section{Frozen Core}

The frozen-core approximation removes a set of low-lying (core) orbitals from
the correlation treatment, splitting the occupied space into core
($c,c',\dots$) and active ($i,j,\dots$) blocks.  The correlation energy remains
invariant to active$\leftrightarrow$active and virtual$\leftrightarrow$virtual
rotations, and to core$\leftrightarrow$core rotations.  It is {\em not}
invariant to core$\leftrightarrow$active-occupied rotations, however, since
mixing a frozen core orbital into the active space changes the set of
correlated excitations.  That block of the MO response therefore requires the
canonical perturbed-MO treatment of \S5 and \S6.

If one chooses to use canonical perturbed MOs (\S5 and \S6), then the extension
of the gradient expressions above to incorporate the MO response remains unchanged.
However, if one is using the non-canonical perturbed MO choice of 
$U^x=-\tfrac12 S^{(x)}$ for the active occ-occ and vir-vir blocks, then one must 
apply the canonical condition \eqref{eq:canon} to the
core$\leftrightarrow$active-occupied block using the corresponding
dependent-pair expression from \S6,
\begin{equation}
P_{ci} = \frac{X_{ci}}{\varepsilon_c-\varepsilon_i}
= \frac{I'_{ci}-I'_{ic}}{\varepsilon_c-\varepsilon_i}
\qquad(c\ \text{core},\ i\ \text{active-occupied}),
\label{eq:Pco}
\end{equation}
the restriction of Eq.~\eqref{eq:canon-oo} to $(c,i)$ pairs.  It enters the
relaxed density as the occupied dependent-pair of \S6 does,
\begin{equation}
\tilde{D}_{ci}=\tilde{D}_{ic}=P_{ci},
\end{equation}
with the overlap-derivative terms folded into $\tilde{I}''$ and the occ-vir-response
coupling into $\tilde{X}$ through Eqs.~\eqref{eq:Idouble-tilde} and
\eqref{eq:Xtilde}, the pair-sum now running over $(c,i)$.  Two further
adjustments distinguish the frozen-core case: (1) 
either the canonical or non-canonical perturbed MO conditions must be applied
to the core-core block; and (2) the occupied$\leftrightarrow$virtual orbital
response (the Z-vector) now spans the {\em full} occupied space, core plus active,
because the core orbitals relax against the virtuals.

Nothing in \S5--\S6 is specific to a particular method or to the all-electron
case.  The occupied dependent-pair is defined over the entire occupied space and
the virtual dependent-pair over the entire virtual space, and each block is
populated only where the energy fails to be invariant, its antisymmetric
numerator $X_{pq}=I'_{pq}-I'_{qp}$ vanishing elsewhere.  

\begin{thebibliography}{1}

\bibitem{Yamaguchi94}
Y. Yamaguchi, Y. Osamura, J.~D. Goddard, and H.~F. Schaefer, \textit{A New
  Dimension to Quantum Chemistry: Analytic Derivative Methods in {\em Ab
  Initio} Molecular Electronic Structure Theory}, No.~29 in
  \textit{International Series of Monographs on Chemistry}, Oxford Univ. Press,
  New York, 1994.

\bibitem{Gauss91:UHF}
J. Gauss, J.~F. Stanton, and R.~J. Bartlett, \textit{J. Chem. Phys.}, {\bf 95},
   2623-2638  (1991).
Coupled-cluster open-shell analytic gradients: Implementation of the direct
  product decomposition approach in energy gradient calculations.

\bibitem{Handy85:MP2Hessians}
N.~C. Handy, R.~D. Amos, J.~F. Gaw, J.~E. Rice, and E.~D. Simandiras,
  \textit{Chem. Phys. Lett.}, {\bf 120},  151  (1985).
The elimination of singularities in derivative calculations.

\bibitem{Handy84}
N.~C. Handy and H.~F. Schaefer, \textit{J. Chem. Phys.}, {\bf 81},  5031-5033
  (1984).
On the evaluation of analytic energy derivatives for correlated wave functions.

\bibitem{Lee91}
T.~J. Lee and A.~P. Rendell, \textit{J. Chem. Phys.}, {\bf 94},  6229-6236
  (1991).
Analytic gradients for coupled-cluster energies that include noniterative
  connected triple excitations: Application to {\em cis-} and {\em trans-}HONO.

\bibitem{Hald03:gradients}
K. Hald, A. Halkier, P. J{\o}rgensen, S. Coriani, and T. Helgaker, \textit{J.
  Chem. Phys.}, {\bf 118},  2985-2998  (2003).
A Lagrangian, integral-density direct formulation and implementation of the
  analytic CCSD and CCSD(T) gradients.

\end{thebibliography}

\end{document}
